(1) 两个格点各有一个自旋轨道，哈密顿量为 \(H=-t\sum_{\sigma=\uparrow,\downarrow}(c_{1\sigma}^\dagger c_{2\sigma}+c_{2\sigma}^\dagger c_{1\sigma})+U\sum_{j=1}^2 n_{j\uparrow}n_{j\downarrow}\)，其中 \(t>0,U>0\)，系统共有两个电子。写出三重态和自旋单态的自然基。
(2) 在自旋单态与对称双占据子空间求 \(2\times2\) 哈密顿量和两个本征能，给出三重态能量。
(3) 求基态单态相对三重态的能量劈裂 \(J\)，并在 \(U\gg t\) 近似下得到超交换的领先结果。
(4) 用 Hellmann–Feynman 定理计算基态总双占据概率，检查 \(U\to0\) 与 \(U\gg t\) 极限并说明 Pauli 原理的作用。
